\section{评估结果}
\label{sec:results}

\subsection{可靠性}
\label{sec:rq-help}
初始研究（1{,}944 个回合，设计者表）比较了完整的 \hesp{}（EIG/成本选择加结束守卫）和 Memory-only；在后者中，规划器看到同样的账本，但自己选探针。Qwen2.5-7B 的已验证完成率提高了 $+0.889$ $[+0.75,+1.00]$，32B 提高了 $+0.250$ $[+0.10,+0.42]$。72B 的提升为 $+0.111$，区间触及零（$[0.000,+0.236]$）。收益随模型规模缩小，与预期一致。这项研究留下一个悬而未决的问题，因为它的表来自环境生成器：收益是否需要神谕知识？表来源研究用计数表替换了设计者表，唯一的主要终点是 7B 与 Memory-only 的比较。\Cref{tab:v06} 报告了每种表来源下的已验证完成率。

\begin{table*}[t]
  \caption{表来源研究：按预测表来源划分的已验证完成率（每格 72 个回合）。7B 与 Memory-only 的比较是预注册的主要终点；32B 和 72B 为探索性结果。区间为 95\% 任务簇自助法区间。表中 empirical 即计数表。}
  \label{tab:v06}
  \centering\small
  \begin{adjustbox}{max width=\textwidth}% generated by paper/make_tables.py -- do not edit
\begin{tabular}{lrrr}
\toprule
Arm & Qwen2.5-7B & Qwen2.5-32B & Qwen2.5-72B \\
\midrule
ReAct-style & 0.028 & 0.611 & 0.597 \\
Memory-only & 0.125 & 0.750 & 0.889 \\
HESP, empirical table ($k{=}20$) & 1.000 & 1.000 & 1.000 \\
HESP, designer table & 1.000 & 1.000 & 1.000 \\
HESP, LLM-elicited table & 0.069 & 0.375 & 0.375 \\
\midrule
$\Delta$ empirical $-$ Memory-only & +0.875 [+0.74, +0.97] & +0.250 [+0.10, +0.42] & +0.111 [+0.00, +0.24] \\
$\Delta$ designer $-$ empirical & +0.000 [+0.00, +0.00] & +0.000 [+0.00, +0.00] & +0.000 [+0.00, +0.00] \\
Mean probe cost, Memory-only / empirical & 9.44 / 9.40 & 5.39 / 6.11 & 5.71 / 5.33 \\
\bottomrule
\end{tabular}
\end{adjustbox}
\end{table*}

有三点发现值得关注。第一，计数表与神谕持平。用每个起因和变体 $k{=}20$ 个开发回合计数得到的表，\hesp{} 把 Qwen2.5-7B 从 0.125 提升到 1.000（$+0.875$ $[+0.74,+0.97]$；22 个任务更好，没有任务更差），并且在三个模型上的完成率都与设计者表相同。$k{=}5$ 时已经饱和，$k{=}1$ 时每个模型最多少一个回合。这种饱和反映的是环境几乎是确定性的，而不是一种普遍的数据效率性质。

第二，不是神谕的代价只集中在一行。计数表比设计者表每回合多花 1.2 到 2.4 个探针单位。一次无 LLM 检查把全部差距定位在\textit{其他}这一行，也就是开发数据中从未出现过的那个起因。只把\textit{其他}这一行换成设计者的，就能精确恢复完成率和成本：在脚本化设置下，从 0.875、3.53 单位变为 1.000、2.54 单位。

第三，模型无法自己写出这些表。使用规划器自报的表时，7B、32B 和 72B 模型的完成率分别降到 0.069、0.375 和 0.375。在结束守卫下，32B 和 72B 模型对撞库、SQL 注入、DNS 命令与控制、暴露的管理面板、有漏洞的组件所引起的案例一个都没完成，而其余三种起因的案例全部完成。因此，控制器必须同时提供流程和表中所编码的知识。

\subsection{排序还是剥夺选择}
\label{sec:rq-rank}
由控制器挑选探针会同时改变两件事：去掉了模型自己的选择，并用一个排序替代了这个选择。分解研究用盲配置把两者分开：无论使用哪个选择器，规划器收到的提示都相同；同时加入 Llama~3.1 家族。它的主要终点是 Qwen2.5-7B 上的排序效应（盲 EIG/成本对盲随机），以及 Llama-3.1-8B 上完整 \hesp{} 对 Memory-only。\Cref{fig:decomposition} 为每个模型展示了收益的两个组成部分；附录~\ref{app:tables} 列出了全部数值。

\begin{figure*}[t]
  \centering
  \includegraphics[width=0.92\textwidth]{fig-decomposition.pdf}
  \caption{分解研究：相对于 Memory-only 的收益来自哪里，逐模型展示（已验证完成率之差，72 个配对回合，任务簇自助法区间）。\emph{剥夺选择}（Taking the choice away）比较控制器随机选择与模型自己的选择。\emph{排序}（Ranking）在规划器提示完全相同时比较 EIG/成本选择与随机选择。空心标记（Qwen 7B，排序）是主要终点，使用 97.5\% 区间；其余为探索性的 95\% 区间。琥珀色标出完全低于零的区间。Llama-3.1-8B 从不下结论，因此它上面的每个对比都为零。}
  \label{fig:decomposition}
\end{figure*}

有三点发现值得关注。第一，排序本身有用。在规划器信息相同的情况下，EIG/成本选择在 Qwen2.5-7B 上比随机选择高 $+0.306$ $[+0.14,+0.47]$，证实了第一个主要终点。在每个会下结论的模型上效果相近：Qwen 32B、Qwen 72B 和 Llama 70B 上分别为 $+0.35$、$+0.35$ 和 $+0.26$，区间都不含零。

第二，剥夺选择对弱模型有帮助，对强模型有害。把模型自己的探针选择换成随机选择，7B 模型的完成率提高 $+0.556$，因为它自己的选择比随机还差；而 Qwen 32B 和 72B 分别降低 $0.222$ 和 $0.278$，因为它们的选择优于随机。因此，对小模型来说，控制器贡献的是纪律和排序；对大模型来说，控制器只贡献排序，而排序之所以能抵消失去的选择，只是因为它信息量足够大。向规划器展示排序、加入结束守卫，各自对完成率的改变都不超过 $+0.06$。

第三，选探针并不能让模型下结论。第二个主要终点未能证实。Llama-3.1-8B 在任何配置下 360 个回合一个都没完成：它做了 3{,}389 次决策，全部是要求再运行一个探针，输出始终是合法 JSON。控制器决定了探什么，但何时停仍由规划器决定。

\subsection{探测还是停止}
\label{sec:rq-stop}
停止研究把“谁选探针”和“谁可以停止”交叉组合，加入随机选择对照，并把两个主要终点都放在 Llama-3.1-8B 上：控制器停止能否挽救它，以及在提供停止之后控制器选探针是否仍有帮助。运行之前，一次无 LLM 模拟预测：如果模型继续从不下结论，第一个终点约为 0.9；因此我们把它视为近乎机械的结果，把第二个终点视为真正有信息量的检验。\Cref{fig:probe-stop} 和 \cref{tab:v09} 报告了结果。

\begin{figure*}[t]
  \centering
  \includegraphics[width=0.92\textwidth]{fig-probe-stop.pdf}
  \caption{停止研究：按“谁选探针”和“谁可以停止”划分的已验证完成率（每个点 72 个回合）。虚线是事后运行的无 LLM 控制器（EIG/成本选择加控制器停止，种子和任务相同）。}
  \label{fig:probe-stop}
\end{figure*}

\begin{table*}[t]
  \caption{停止研究：已验证完成率，括号内为平均探针成本，每格 72 个回合。“LLM or controller”：规划器仍可抢先结束。新的种子（2027）使第一行和第三行成为分解研究的独立重复。底部几行用一个从不结束的脚本替换了 LLM。}
  \label{tab:v09}
  \centering\small
  \begin{adjustbox}{max width=\textwidth}% generated by paper/make_tables.py -- do not edit
\begin{tabular}{llrrrrr}
\toprule
Who probes & Who stops & Q-7B & Q-32B & Q-72B & L-8B & L-70B \\
\midrule
LLM & LLM & 0.083 {\scriptsize(9.5)} & 0.903 {\scriptsize(5.2)} & 0.861 {\scriptsize(5.6)} & 0.000 {\scriptsize(9.7)} & 0.889 {\scriptsize(8.6)} \\
LLM & LLM or controller & 0.694 {\scriptsize(5.2)} & 0.819 {\scriptsize(4.1)} & 0.833 {\scriptsize(4.1)} & 0.861 {\scriptsize(4.8)} & 0.889 {\scriptsize(4.4)} \\
controller (EIG/cost) & LLM & 0.931 {\scriptsize(7.2)} & 1.000 {\scriptsize(5.5)} & 0.917 {\scriptsize(3.6)} & 0.000 {\scriptsize(10.0)} & 1.000 {\scriptsize(9.0)} \\
controller (EIG/cost) & LLM or controller & 0.917 {\scriptsize(2.3)} & 0.917 {\scriptsize(2.3)} & 0.875 {\scriptsize(2.3)} & 0.917 {\scriptsize(2.3)} & 0.917 {\scriptsize(2.3)} \\
controller (random) & LLM or controller & 0.736 {\scriptsize(6.4)} & 0.708 {\scriptsize(6.0)} & 0.694 {\scriptsize(5.4)} & 0.750 {\scriptsize(6.4)} & 0.750 {\scriptsize(6.4)} \\
\midrule
\multicolumn{7}{l}{\textit{No LLM (post hoc reference, same seeds, tasks and budget)}} \\
catalogue order & controller & \multicolumn{5}{c}{0.847 {\scriptsize(4.8)}} \\
controller (EIG/cost) & controller & \multicolumn{5}{c}{0.917 {\scriptsize(2.3)}} \\
controller (random) & controller & \multicolumn{5}{c}{0.750 {\scriptsize(6.4)}} \\
\bottomrule
\end{tabular}
\end{adjustbox}
\end{table*}

第一，控制器停止挽救了从不下结论的模型。允许控制器停止后，Llama-3.1-8B 从 0 升到 0.917（$+0.917$ $[+0.79,+1.00]$；22 个任务更好，没有任务更差），与预测一致。

第二，对这个模型而言，收益来自停止。提供停止之后，Llama-3.1-8B 自己选探针达到 0.861，控制器选探针达到 0.917，差值为 $+0.056$ $[-0.07,+0.18]$。该区间允许存在中等程度的收益，因此我们只能得出：对这个模型，选择没有带来可检测的完成率提升。但选择确实把探针成本减半（2.35 对 4.78 单位）。模型自己的选择表现与固定的目录顺序相当，后者在无 LLM 运行中以 4.85 单位达到 0.847。

第三，小模型的失败方式不同。仅靠控制器停止就把 Qwen2.5-7B 从 0.083 提升到 0.694，控制器选探针再增加 $+0.222$ $[+0.08,+0.39]$（附录~\ref{app:tables}）。因此它在分解研究中的失败，部分是拒绝下结论，部分是探测不佳；而 Llama-3.1-8B 的失败几乎完全是拒绝下结论。控制器无法事先知道某个模型缺的是哪一部分，所以两部分都得提供。

第四，排序效应在全部五个模型上重复出现。在控制器停止下，EIG/成本在每个模型上都比随机选择高 $+0.17$ 到 $+0.21$，包括 Llama-3.1-8B，区间都不含零。

第五，控制器靠排除法下结论，而强模型会等待确认。在本来就会下结论的模型上，允许控制器停止会使完成率降低 0.01 到 0.08。控制器停止配置的每一个失败（每个模型六个回合）都是一个正确的 DNS 命令与控制结论，由控制器排除其他可能后得出。验证器要求确认性特征，因此拒绝了它，这在运行前的协议中就已预料到。当由 Qwen 32B 和 Llama 70B 决定何时停止时，它们会继续探测直到 DNS 证据出现，并完成全部案例。唯一的其他失败是三个 drift 回合，其中 Qwen 72B 抢先以错误的起因结束。

最后，控制器本身已接近最好水平。无 LLM 运行以 2.35 单位完成 0.917 的案例，是所有配置中成本最低的。当控制器既选探针又负责停止时，五个模型中有四个也恰好得到 0.917、成本相同，说明模型在这里几乎没有贡献。超过它的配置，都是由强模型决定证据何时已足够的配置。

\subsection{失败方式}
\label{sec:security}
\Cref{tab:security} 按错误类型拆分了表来源研究中的结论。这些指标是描述性的，没有预注册为终点。

\begin{table*}[t]
  \caption{表来源研究：面向安全的结果（描述性）。“claimed”统计以结论结束的回合。漏报攻击（missed attack）、误升级（false escalation）和错因（wrong cause）率以这些结论为分母；未解决（unresolved）率以全部 72 个回合为分母。}
  \label{tab:security}
  \centering\small
  \begin{adjustbox}{max width=\textwidth}% generated by paper/make_tables.py -- do not edit
\begin{tabular}{llrrrrr}
\toprule
Model & Arm & claimed & unresolved & missed attack & false escalation & wrong cause \\
\midrule
Qwen2.5-7B & ReAct-style & 3/72 & 0.97 & 1.00 & 0.00 & 0.33 \\
Qwen2.5-7B & Memory-only & 9/72 & 0.88 & 0.00 & 0.00 & 0.00 \\
Qwen2.5-7B & HESP (empirical) & 72/72 & 0.00 & 0.00 & 0.00 & 0.00 \\
\midrule
Qwen2.5-32B & ReAct-style & 69/72 & 0.39 & 0.25 & 0.28 & 0.36 \\
Qwen2.5-32B & Memory-only & 72/72 & 0.25 & 0.04 & 0.17 & 0.12 \\
Qwen2.5-32B & HESP (empirical) & 72/72 & 0.00 & 0.00 & 0.00 & 0.00 \\
\midrule
Qwen2.5-72B & ReAct-style & 72/72 & 0.40 & 0.37 & 0.17 & 0.40 \\
Qwen2.5-72B & Memory-only & 71/72 & 0.11 & 0.02 & 0.00 & 0.06 \\
Qwen2.5-72B & HESP (empirical) & 72/72 & 0.00 & 0.00 & 0.00 & 0.00 \\
\bottomrule
\end{tabular}
\end{adjustbox}
\end{table*}

有三点观察值得关注。第一，使用计数表时，\hesp{} 在全部 216 个回合中都给出了结论，没有一个给错起因。第二，基线的失败方式随规模而变。7B 基线很少下结论，88\% 到 97\% 的回合未解决；而更大的 ReAct 式基线会下结论但经常出错：它们把 25\%（32B）和 37\%（72B）的需处置事件判为良性，引用中分别有 22\% 和 31\% 是无效证据。第三，仅账本本身就消除了大部分这类错误：Memory-only 把漏报攻击率降到 4\% 和 2\%。在停止研究中，控制器选探针加控制器停止的配置在全部五个模型上都从未把需处置事件判为良性。

\subsection{公开安全数据}
\label{sec:realdata}
受控环境引出一个显而易见的问题：这些结果在真实数据上还成立吗？我们尝试用最接近的三个公开数据集来回答。下面的测量只使用训练或开发材料，并遵循同样的预注册纪律。不过，两个更早的可行性检查脚本读取过 ExCyTIn 全部 1{,}017 道题，其中包括 599 道测试题（勘误 E-8）；没有任何 ExCyTIn 结果用于评估 \hesp{}，GUIDE 的测试集也始终未被读取。要检验控制器是否选对探针，需要满足三点。（i）仅凭告警本身不能揭示答案，因此必须有多种解释仍然可能。（ii）探针的成本和揭示的信息量必须不同，这样顺序才有意义。（iii）真值必须客观，反映实际发生了什么，而不是一种标注惯例。三个数据集都无法同时满足这三点。

\paragraph{ExCyTIn-Bench~\citep{excytin}} ExCyTIn 针对一个模拟 Azure 租户的日志提出问题，其答案可以从公开的调查图中机械地恢复（1{,}017 个问题中的 98.5\%）。在 418 个训练问题中，54.8\% 在问题文本里直接点名了目标告警，其余大多数是对它的转述，因此任务变成定位被点名的告警并读出一个实体。标准的指标关联探针找到真实答案的频率并不比干扰项高（告警表中为 0.571 对 0.570），只有 19.4\% 的问题存在能单独挑出答案的表。这些问题检验的是导航和检索，而不是调查。

\paragraph{GUIDE~\citep{guide}} GUIDE 包含来自 5{,}340 个组织的 707{,}108 起真实、已匿名的事件，标为真阳性、良性阳性或假阳性。在我们基于其告警元数据定义的十个事件级探针上，组织身份本身就携带了等级 1.50 比特熵中的 1.03 比特，还有几个大型组织对每起事件都给同一个等级。在组织内部，查一下该组织以前如何给同一检测器定级，就能达到 0.911 的准确率；跨组织时，这些探针不超过多数类基线。在冷启动事件上，证据整合（AUROC 0.755）与查历史（0.759）持平（差值 $-0.004$ $[-0.028,+0.027]$）。标签反映的是组织策略，因此我们暂停了计划中的 GUIDE 研究；它的测试集至今未读。

\paragraph{OTRF Security-Datasets~\citep{otrf}} OTRF 包含执行 ATT\&CK 技术时在实验室中录制的 Windows 主机遥测，没有良性活动的录制。我们审计了远程执行类别的 27 份录制（243{,}000 个事件），把每次服务安装和计划任务的创建或更新当作一条告警，并用录制自带的元数据标注。在 46 条不同的告警中，8 条是攻击，38 条是后台操作系统事件。一条只看告警文本的简单规则——判断该任务或服务是否属于 Windows 组件——在 46 条中判对 44 条，因为 8 次攻击中有 6 次被其工具的名字暴露了。

\paragraph{案例研究} OTRF 中唯一确实需要调查的场景，在真实日志上展示了 \hesp{} 所针对的那类案例。它的元数据引用了微软对 Solorigate 入侵的分析~\citep{solorigate}，并把该活动标为远程创建计划任务（T1053.005）。告警是 \texttt{WORKSTATION6} 上的安全事件 4698，报告在文件夹 \texttt{\textbackslash Microsoft\textbackslash Windows\textbackslash Software\-Protection\-Platform} 中新建了任务 \texttt{EventCacheManager}，其文件夹和名字都与真实的 Windows 任务一致；只看文本的规则判定它为良性。第一个探针发现它是域用户 \texttt{THESHIRE\textbackslash pgustavo} 创建的，而 27 份录制中 \texttt{\textbackslash Microsoft\textbackslash Windows} 下的另外 17 个任务事件全部由机器账户创建。第二个探针发现创建它的会话是来自另一台主机的网络登录（类型 3，Kerberos），这排除了本地持久化和用户上下文更新程序。该任务以 SYSTEM 身份运行 \texttt{cmd.exe}，源主机的进程日志显示了创建它的 \texttt{schtasks /create \dots{} /s} 命令。告警文本和调查结果指向相反的方向，两个便宜的探针就能把它们区分开。由于一个探针就足够、而且场景只有这一个，这个案例无法说明探针顺序是否重要；我们只把它作为定性证据。

\paragraph{共同的缺口} 三个数据集的失败方式各不相同，但缺的是同一种东西：看起来像日常运维的攻击，以及看起来像攻击的日常运维。它们的良性样本是操作系统的日常维护，而不是管理员合法使用同样的工具；它们的攻击模拟留下了谨慎的攻击者不会留下的痕迹。这类成对的、客观标注的数据，是在真实遥测上评估分诊智能体的前提。
